%%%%%%%%%%%%%%%%%%%%%%%%%%%%%%%%%%%%%%%%%%%%%%%%
% Python Cheat Sheet
% baposter Landscape Poster
% LaTeX Template
% Version 1.0 (11/06/13)
% baposter Class Created by:
% Brian Amberg (baposter@brian-amberg.de)
% This template has been downloaded from:
% http://www.LaTeXTemplates.com
% License:
% CC BY-NC-SA 3.0 (http://creativecommons.org/licenses/by-nc-sa/3.0/)
% Edited by Michelle Cristina de Sousa Baltazar
%%%%%%%%%%%%%%%%%%%%%%%%%%%%%%%%%%%%%%%%%%%%%%%%

%----------------------------------------------------------------
%	PACKAGES AND OTHER DOCUMENT CONFIGURATIONS
%----------------------------------------------------------------

\documentclass[landscape,a0paper,fontscale=0.333]{baposter} % Adjust the font scale/size here
% \title{Python Cheat Sheet New}
\usepackage[utf8]{inputenc}

% \renewcommand{\familydefault}{<family>}
\renewcommand{\familydefault}{\sfdefault}

\usepackage{graphicx} % Required for including images
\graphicspath{{figures/}} % Directory in which figures are stored

\usepackage{xcolor}
\usepackage{colortbl}
\usepackage{tabu}

\usepackage{mathtools}
%\usepackage{amsmath} % For typesetting math
\usepackage{amssymb} % Adds new symbols to be used in math mode

\usepackage{booktabs} % Top and bottom rules for tables
\usepackage{enumitem} % Used to reduce itemize/enumerate spacing
\usepackage{palatino} % Use the Palatino font
\usepackage[font=small,labelfont=bf]{caption} % Required for specifying captions to tables and figures

\usepackage{multicol} % Required for multiple columns
\setlength{\columnsep}{1.5em} % Slightly increase the space between columns
\setlength{\columnseprule}{0mm} % No horizontal rule between columns

\usepackage{algorithm}
\usepackage[noend]{algpseudocode}

\usepackage{tikz} % Required for flow chart
\usetikzlibrary{decorations.pathmorphing}
\usetikzlibrary{shapes,arrows} % Tikz libraries required for the flow chart in the template

\newcommand{\compresslist}{%
    % Define a command to reduce spacing within itemize/enumerate
    % environments, this is used right after \begin{itemize}
    % or \begin{enumerate}
    \setlength{\itemsep}{1pt}
    \setlength{\parskip}{0pt}
    \setlength{\parsep}{0pt}
}


\usepackage{amsthm}
\newtheorem{thm}{Theorem}
\newtheorem{cor}[thm]{Corollary}
\newtheorem{lemma}[thm]{Lemma}
\theoremstyle{definition}
\newtheorem{defn}{Definition}
\newtheorem{remark}{Remark}

\usepackage[usenames,dvipsnames]{xcolor}
\usepackage{hyperref}
\hypersetup{
    unicode=false,                % non-Latin characters in Acrobat’s bookmarks
    pdftoolbar=true,              % show Acrobat’s toolbar?
    pdfmenubar=true,              % show Acrobat’s menu?
    pdffitwindow=false,           % window fit to page when opened
    pdfstartview={FitH},          % fits the width of the page to the window
    pdftitle={CV},                % title
    pdfauthor={Scott Sievert},    % author
    pdfsubject={},                % subject of the document
    pdfcreator={Scott Sievert},   % creator of the document
    pdfproducer={Scott Sievert},  % producer of the document
    pdfkeywords={compressed sensing,} {research,} {adaptive sampling,}, % list of keywords
    pdfnewwindow=true,            % links in new window
    colorlinks=true,              % false: boxed links; true: colored links
    linkcolor=red,                % color of internal links
    citecolor=green,              % color of links to bibliography
    filecolor=magenta,            % color of file links
    urlcolor=NavyBlue             % color of external links
}

\newcommand{\yb}{\mathbf{y}}
\newcommand{\wb}{\mathbf{w}}
\newcommand{\bb}{\mathbf{b}}
\newcommand{\xb}{\mathbf{x}}
\newcommand{\xhatb}{\widehat{\mathbf{x}}}
\newcommand{\Ab}{\mathbf{A}}
\newcommand{\Xb}{\mathbf{X}}
\newcommand{\matrix}[1]{\begin{bmatrix} #1 \end{bmatrix}}
\newcommand{\norm}[1]{\left\lVert #1 \right\rVert}
\newcommand{\abs}[1]{\left| #1 \right|}
\newcommand{\Align}[1]{\begin{aligned} #1 \end{aligned}}
\newcommand{\parens}[1]{\left( #1 \right)}
\newcommand{\braces}[1]{\left\{ #1 \right\}}
\newcommand{\brackets}[1]{\left #1 \right]}
\newcommand{\deriv}{\frac{d}{dx}}
\newcommand{\prob}[1]{\Pr\parens{#1}}
\newcommand{\at}[1]{\Big|_{#1}}
\newcommand{\reals}{\mathbb{R}}
\newcommand{\R}{\mathbb{R}}
\newcommand{\given}[1][]{\:\Big\vert\:}


\usepackage{xcolor}
\definecolor{headerbg}{RGB}{220,220,255} % Defines the color used for content box headers
%\definecolor{headerbg}{RGB}{255,220,220} % Defines the color used for content box headers
\definecolor{subheaderbg}{RGB}{230,230,255} % Defines the color used for content box headers
\definecolor{gray}{RGB}{200,200,200} % Defines the color used for content box headers
\newcommand{\subheading}[1]{%
    \par\noindent\colorbox{subheaderbg}
        {\parbox{\dimexpr\textwidth-2\fboxsep\relax}{\textbf{\large #1}}}
}


\begin{document}

\begin{poster}{
    headerborder=closed, % Adds a border around the header of content boxes
    colspacing=0.8em, % Column spacing
    bgColorOne=white, % Background color for the gradient on the left side of the poster
    bgColorTwo=white, % Background color for the gradient on the right side of the poster
    borderColor=gray, % Border color
    headerColorOne=headerbg, % Background color for the header in the content boxes (left side)
    headerColorTwo=headerbg, % Background color for the header in the content boxes (right side)
    headerFontColor=black, % Text color for the header text in the content boxes
    boxColorOne=white, % Background color of the content boxes
    textborder=roundedleft, % Format of the border around content boxes, can be: none, bars, coils, triangles, rectangle, rounded, roundedsmall, roundedright or faded
    eyecatcher=true, % Set to false for ignoring the left logo in the title and move the title left
    headerheight=0.1\textheight, % Height of the header
    headershape=roundedright, % Specify the rounded corner in the content box headers, can be: rectangle, small-rounded, roundedright, roundedleft or rounded
    headerfont=\Large\bf\textbf, % Large, bold and sans serif font in the headers of content boxes
    %textfont={\setlength{\parindent}{1.5em}}, % Uncomment for paragraph indentation
    linewidth=2pt, % Width of the border lines around content boxes
    columns=3,
}
{\includegraphics[width=0.15\textwidth]{uw-madison-logo-wide}}
{Foo \ \ Bar \ \ for \ \ a \ \ ERM \ \ with \ \ fancy \ \ method}
{\LARGE Foo \  Bar and \  Goo \  Baz\\UW--Madison ECE}
{\includegraphics[width=0.15\textwidth]{WIDBlackHorizontalPlain-1}}


% First column
\headerbox{Abstract}{name=c0r0,column=0,row=0,span=1}{%
    Te mea choro apeirian adversarium. Ei insolens oportere mel, ei agam qualisque
per. Usu id placerat dissentiunt, ea pri verear assueverit. Nulla saepe
assueverit te vim, in vel animal antiopam expetendis, apeirian salutandi vix
ad. Ex sed stet minimum, eam elitr intellegebat cu, te eros augue nostro sit.
ad. Ex sed stet minimum, eam elitr intellegebat cu, te eros augue nostro sit.
Nam accumsan mentitum hendrerit no, noster utroque pri eu, at his possit
referrentur efficiantur.
Nam accumsan mentitum hendrerit no, noster utroque pri eu, at his possit
referrentur efficiantur.

}

\headerbox{Background}{name=c0r1,column=0,row=1,below=c0r0, span=1}{%
    Aperiam evertitur has in. Cu brute scaevola dis

$$\min_{\xb\in\R^d}F(\xb)$$

que efficiendi mel, eos $f_i$ cu facilisis $i$ intellegam. Vix quidam invidunt et, ex mea amet alia. Impetus maioru:

$$\xb_{k+1} = \Pi\left(\xb_k - \gamma\grad F(\xb)\right)$$

ribentur ad mei, dolore viderer epicurei at nam, admodum copiosae his id. Ocurreret elaboraret interesset pro eu. Eirmod laboramus assentior in vim, percipit hendrerit moderatius ex sed.

}
\headerbox{Problem}{name=c0r2,column=0,row=1,below=c0r1, span=1}{%
    m id eos, oratio melius quaerendum has ea. Minim veritus an

\begin{itemize}\compresslist
    \item ea. Minim veritus an
    \item va ase Minim veritus anea.
\end{itemize}

quo, facer populo nostrud ex has. Maluisset hendrerit his no, ius choro perpetua interpretaris cu.\\

 Vivendo inciderint id his. Qui modus consul delicatissimi an, primis

\begin{itemize}
        \compresslist
    \item accommodare ad, sed praesent mnesarchum ei. Euismod abhorreant pro at, te mutat regione pericula mel, id est mundi veniam inimicus.
    \item eque urbanitas vel, ex eos atqui iriure ceteros, sit diam assum alienum an. Ius te sale \textit{appetere} ap
\end{itemize}

Petentium vituperata cu sea. Alienum verterem ex pro $F(\xb)$. Tantas offendit eum ad. Accusata intellegam quo cu:

\hspace{5em}\includegraphics[width=0.60\textwidth]{graph}

mod interpretaris. Nec facete iudicabit liberavisse ex, lorem reformidans consectetuer ut v
\cite{bubeck2015convex}, and increases.

}

% Second column
\headerbox{Algorithm}{name=c1r0,column=1,row=0,span=1}{%
    
ed utinam explicari an, virtute theophrastus nam ea. Enim everti legimus quo ex, ius ea possit minimum:
$$\gamma_k = \gamma_0 / \sqrt{k}$$
for some constant $\gamma_0$.

\begin{algorithm}[H]
\caption{Step size decay SGD}\label{alg:adaptive-max-batch}
\begin{algorithmic}[1]
\Procedure{DecaySGD}{step size $\gamma_0$, initial model $\xb_0$, batch size $B$}
    \For{$k\in [0, 1, 2, \ldots]$}
    \If{$k = 0$}
    \State $\gamma_k = \gamma_0$
    \EndIf
    \State $\gamma_k \assign \gamma_0 / \sqrt{k}$
    \State $\est{g} \assign \frac{1}{B}\sum_{j=1}^{B} \grad f_{i_j}(\xb)$
    \State $\xb_{k+1} \assign \xb_k - \gamma_k\est{g}$
    \EndFor
    \Return{$x_k$}
\EndProcedure
\end{algorithmic}
\end{algorithm}

An aeque urbanitas vel, ex eos atqui iriure ceteros, sit diam assum alienum an. Ius te sale appetere appellantur. Ne eum prompta utroque, et verterem argumentum mediocritatem cum. Sit summo mediocritatem te, ei mel.

}
\headerbox{Convergence}{name=c1r1,column=1,row=1,below=c1r0,span=1}{%
    
\subheading{Model updates}

\begin{thm}(informal)
    When algorithm~\ref{alg:adaptive-max-batch} is used,
    When algorithm~\ref{alg:adaptive-max-batch} is used,
    When algorithm~\ref{alg:adaptive-max-batch} is used,
    $$
    e = mc^2
$$
sit summo mediocritatem $c$ and $m$ te, ei mel qualisque interpretaris.
\end{thm}

This obtains a similar solution \cite[Theorem 1.1]{bubeck2015convex}.\\

An aeque urbanAn aeque urbanitas vel, ex eos atqui iriure ceteros, sit diam assum alienum an. itas vel, ex eos atqui iriure ceteros, sit diam assum alienum an. Ius te sale appetere appellantur.\\

\subheading{Number of examples}
By some measures, the number of model updates is meaningless. How many examples
does this optimization method require?

\begin{thm}
    Algorithm~\ref{alg:adaptive-max-batch} requires no more $c = 299 792 458$ m/s to obtain
    a model with loss $F(\xb_T) - F^\star \le \epsilon$ when
    the world has a finite duration and if money isn't free.
\end{thm}

CVX with $\gamma_0 = 1$ also requires $c = 299792458$ m/s \cite{bubeck2015convex}, as does $c' = c/2$.

}

% Third column
\headerbox{Experiments}{name=c2r0,column=2,row=0}{%
    \input{src/experiments}
}
\headerbox{Conclusion}{name=c2r1,column=2,row=1,below=c2r0}{%
    \input{src/conclusion}
\bibliography{src/refs}{}
\bibliographystyle{plain}
}
\end{poster}
\end{document}
